\begin{tikzpicture}[font=\small,>=stealth,line join=round]
\draw[gray!40,rounded corners=3pt] (0,0) rectangle (14.2,7.8);
\node at (7.1,7.2) {去分母后得到的根，可能恰好落在原函数的孔点};
\begin{scope}[shift={(4.4,2.0)},x=1cm,y=2cm]
\draw[->,gray] (-3.4,0)--(3.4,0) node[right]{$x$};\draw[->,gray] (0,-.2)--(0,2.2) node[above]{$y$};
\draw[blue,very thick] (-3.1,1)--(3.1,1);
\draw[orange!85!black,thick,fill=white] (-1,1) circle(3pt);\draw[orange!85!black,thick,fill=white] (1,1) circle(3pt);
\node[above left=7pt] at (-1,1){$(-1,1)$};\node[above right=7pt] at (1,1){$(1,1)$};
\node[below left=4pt] at (0,0){$O$};
\end{scope}
\node[align=left,anchor=west] at (8.4,4.2){$y=\dfrac{x^2-1}{x^2-1}$\\[12pt]$x\ne\pm1$\\[12pt]两处孔点不能当作解\\[8pt]真正值域仍为 $\{1\}$};
\node[text width=13cm,align=center] at (7.1,.6) {空心点对应分母为零的位置；它们不能证明其他 $y$ 可以取到。};
\end{tikzpicture}
